\begin{filecontents*}{references.bib}

@article{Harkonen1987,
  author  = {H{\"a}rk{\"o}nen, T.},
  title   = {Seasonal and regional variations in the feeding habits of the harbour seal, \textit{Phoca vitulina}, in the {S}kagerrak and {K}attegat},
  journal = {Journal of Zoology},
  year    = {1987},
  volume  = {213},
  pages   = {535--543}
}

@article{Harkonen1988a,
  author  = {H{\"a}rk{\"o}nen, T.},
  title   = {Food habits of the harbour seal (\textit{Phoca vitulina}) in the {S}kagerrak},
  journal = {Journal of Zoology},
  year    = {1988},
  volume  = {214},
  pages   = {673--681}
}

@phdthesis{Harkonen1988b,
  author  = {H{\"a}rk{\"o}nen, T.},
  title   = {Ecology of the harbour seal (\textit{Phoca vitulina}) in the {S}kagerrak--{K}attegat},
  school  = {University of Gothenburg},
  year    = {1988}
}

@article{HarkonenHeide1991,
  author  = {H{\"a}rk{\"o}nen, T. and Heide-J{\o}rgensen, M. P.},
  title   = {The harbour seal \textit{Phoca vitulina} as a predator in the {S}kagerrak},
  journal = {Ophelia},
  year    = {1991},
  volume  = {34},
  pages   = {191--207}
}

@incollection{OlsenBjorge1995,
  author    = {Olsen, M. T. and Bj{\o}rge, A.},
  title     = {Seasonal and regional variation in the diet of the harbour seal in {N}orwegian waters},
  booktitle = {Developments in Marine Biology; Whales, Seals, Fish and Man},
  volume    = {4},
  pages     = {271--285},
  year      = {1995}
}

@techreport{Stromberg2012,
  author      = {Str{\"o}mberg, A. and others},
  title       = {Harbour seal diet in {S}wedish waters},
  institution = {Swedish Museum of Natural History},
  year        = {2012},
  number      = {Report No.~5: 2012}
}

@article{Sorlie2020,
  author  = {S{\o}rlie, M. A. and others},
  title   = {Diet of harbour seals (\textit{Phoca vitulina}) along the {N}orwegian coast},
  journal = {Marine Biology Research},
  year    = {2020},
  volume  = {16},
  pages   = {299--310}
}

@article{ScharffOlsen2019,
  author  = {Scharff-Olsen, C. H. and others},
  title   = {Diet of seals in the {B}altic {S}ea region: a synthesis of published and new data from 1968 to 2013},
  journal = {ICES Journal of Marine Science},
  year    = {2019},
  volume  = {76},
  pages   = {284--297},
  doi     = {10.1093/icesjms/fsy159}
}

@article{BowenIverson2013,
  author  = {Bowen, W. D. and Iverson, S. J.},
  title   = {Methods of estimating marine mammal diets: a review of validation experiments and sources of bias and uncertainty},
  journal = {Marine Mammal Science},
  year    = {2013},
  volume  = {29},
  pages   = {719--754}
}

@article{Pitcher1980,
  author  = {Pitcher, K. W.},
  title   = {Stomach contents and feces as indicators of harbor seal, \textit{Phoca vitulina}, foods in the {G}ulf of {A}laska},
  journal = {Fishery Bulletin},
  year    = {1980},
  volume  = {78},
  pages   = {797--798}
}

@article{Young2010,
  author  = {Young, B. G. and others},
  title   = {Variation in the diet of harbour seals (\textit{Phoca vitulina}) in response to changes in prey availability},
  journal = {Polar Biology},
  year    = {2010},
  volume  = {33},
  number  = {2},
  pages   = {153--162},
  doi     = {10.1007/s00300-009-0693-3}
}

@techreport{WilsonHammond2016,
  author      = {Wilson, L. J. and Hammond, P. S.},
  title       = {The diet of grey and harbour seals around {U}{K}: Examining the role of prey availability and implications for population dynamics},
  institution = {Scottish Marine and Freshwater Science},
  year        = {2016},
  number      = {7(21)}
}

@article{Neises2021,
  author  = {Neises, V. and others},
  title   = {Age-related dietary variation in harbour seals (\textit{Phoca vitulina}) along the {G}erman coast},
  journal = {Conservation Physiology},
  year    = {2021},
  volume  = {9},
  number  = {1},
  doi     = {10.1093/conphys/coab036}
}

@article{Conwell2024,
  author  = {Conwell, C. and others},
  title   = {Sex and age differences in diet composition of {G}rey seals (\textit{Halichoerus grypus})},
  journal = {Ecology and Evolution},
  year    = {2024},
  volume  = {14},
  number  = {7},
  pages   = {e11417},
  doi     = {10.1002/ece3.11417}
}

@article{HansenHarding2006,
  author  = {Hansen, M. and Harding, K. C.},
  title   = {Harbour seal diet in {S}candinavian waters},
  journal = {Journal of Sea Research},
  year    = {2006},
  volume  = {56},
  number  = {4},
  pages   = {329--337}
}

@article{MionOtolith2024,
  author  = {Mion, M. and others},
  title   = {Species assignment from seal diet samples using shape analyses in a machine learning framework},
  journal = {ICES Journal of Marine Science},
  year    = {2024},
  doi     = {10.1093/icesjms/fsae134}
}

@article{Anderson2024,
  author  = {Anderson, S. C. and others},
  title   = {{sdmTMB}: an {R} package for fast, flexible, and user-friendly generalized linear mixed effects models with optional spatial and spatiotemporal random fields},
  journal = {bioRxiv},
  year    = {2024},
  doi     = {10.1101/2022.03.24.485545}
}

@article{LindmarkCod2024,
  author  = {Lindmark, M. and others},
  title   = {Quantifying competition between two demersal fish species from spatiotemporal stomach content data},
  journal = {bioRxiv},
  year    = {2024},
  doi     = {10.1101/2024.04.22.590538}
}

@article{HeideHarkonen1988,
  author  = {Heide-J{\o}rgensen, M. P. and H{\"a}rk{\"o}nen, T.},
  title   = {Rebuilding seal stocks in the {K}attegat-{S}kagerrak},
  journal = {Marine Mammal Science},
  year    = {1988},
  volume  = {4},
  number  = {3},
  pages   = {231--246}
}

@article{Silva2021,
  author  = {Silva, W. T. A. F. and others},
  title   = {Risk for overexploiting a seemingly stable seal population: influence of multiple stressors and hunting},
  journal = {Ecosphere},
  year    = {2021},
  volume  = {12},
  number  = {1},
  pages   = {e03343},
  doi     = {10.1002/ecs2.3343}
}

@article{ICES2023,
  author  = {{ICES}},
  title   = {Report of the {ICES} Working Group on {M}arine {M}ammal {E}cology ({WGMME})},
  journal = {ICES Scientific Reports},
  year    = {2023},
  volume  = {5},
  number  = {88},
  doi     = {10.17895/ices.pub.24131736}
}

@techreport{Larsson2024,
  author      = {Larsson, P. and others},
  title       = {Knubbsälspopulationens i {S}verige},
  institution = {Sveriges lantbruksuniversitet},
  year        = {2024},
  number      = {Aqua notes 2024:6},
  doi         = {10.54612/a.1p05qqkdqa}
}

@article{Cardinale2017,
  author  = {Cardinale, M. and others},
  title   = {Rebuilding depleted fish stocks: the actual rationale of the {EU} {C}ommon {F}isheries {P}olicy},
  journal = {Marine Policy},
  year    = {2017},
  volume  = {83},
  pages   = {179--183}
}

@article{Carroll2025,
  author  = {Carroll, D. and Ahola, M. P. and Carlsson, A. M. and Galatius, A. and Nilssen, K. T. and H{\"a}rk{\"o}nen, T. and Harding, K. C.},
  title   = {Declining harbour seal abundance in a previously recovering meta-population},
  journal = {PLOS ONE},
  year    = {2025},
  volume  = {20},
  number  = {6},
  pages   = {e0326933},
  doi     = {10.1371/journal.pone.0326933}
}

@article{Hilborn2017,
  author  = {Hilborn, R. and others},
  title   = {When does fishing forage species affect their predators?},
  journal = {Fisheries Research},
  year    = {2017},
  volume  = {191},
  pages   = {211--221}
}

@article{BowenLidgard2013,
  author  = {Bowen, W. D. and Lidgard, D.},
  title   = {Marine mammal culling programs: review of effects on predator and prey populations},
  journal = {Mammal Review},
  year    = {2013},
  volume  = {43},
  pages   = {207--220}
}

@article{daSilva1985,
  author  = {da Silva, J. and Neilson, J. D.},
  title   = {Limitations of using otoliths recovered in scats to estimate prey consumption in seals},
  journal = {Canadian Journal of Fisheries and Aquatic Sciences},
  year    = {1985},
  volume  = {42},
  number  = {8},
  pages   = {1439--1442},
  doi     = {10.1139/f85-180}
}

@article{GrEllier2005,
  author  = {Grellier, K. and Hammond, P. S.},
  title   = {Feeding method affects otolith digestion in captive grey seals: Implications for diet composition estimation},
  journal = {Marine Mammal Science},
  year    = {2005},
  volume  = {21},
  pages   = {296--306}
}

@article{Tverin2019,
  author  = {Tverin, M. and others},
  title   = {Complementary methods assessing short and long-term prey of a marine top predator -- Application to the grey seal--fishery conflict in the {B}altic {S}ea},
  journal = {PLOS ONE},
  year    = {2019},
  volume  = {14},
  number  = {1},
  doi     = {10.1371/journal.pone.0208229}
}

@article{Lundstrom2010,
  author  = {Lundstr{\"o}m, K. and others},
  title   = {Understanding the diet composition of marine mammals: grey seals (\textit{Halichoerus grypus}) in the {B}altic {S}ea},
  journal = {ICES Journal of Marine Science},
  year    = {2010},
  volume  = {67},
  number  = {6},
  doi     = {10.1093/icesjms/fsq022}
}

@article{Harding2018,
  author  = {Harding, K. C. and others},
  title   = {Evaluating the sustainability of wildlife harvest: uncertainty, biases, and trade-offs in pinniped hunting},
  journal = {Frontiers in Ecology and Evolution},
  year    = {2018},
  volume  = {6},
  pages   = {e59},
  doi     = {10.3389/fevo.2018.00059}
}

\end{filecontents*}

\documentclass[12pt]{article}
\usepackage[utf8]{inputenc}
\usepackage[T1]{fontenc}
\usepackage{amsmath}
\usepackage{graphicx}
\usepackage[round,authoryear]{natbib}
\usepackage[hidelinks]{hyperref}
\usepackage{lineno}
\usepackage{booktabs}
\usepackage{caption}
\usepackage[margin=2.5cm]{geometry}
\usepackage{setspace}
\usepackage{microtype}
\usepackage{xcolor}
\usepackage{array}
\bibliographystyle{apalike}

\doublespacing

\title{Body size, sex, and season shape the diet of harbour seals (\textit{Phoca vitulina}) in the Skagerrak: evidence from spatiotemporal modelling of hunted individuals}

\author{
  Monica Mion$^{1,*}$ \and
  Max Lindmark$^{1}$ \and
  Karl Lundstr{\"o}m$^{1}$ \and
  Karin H{\aa}rding$^{2}$
}

\date{
  $^{1}$Department of Aquatic Resources, Swedish University of Agricultural Sciences (SLU), Lysekil, Sweden \\
  $^{2}$Department of Biological and Environmental Sciences, University of Gothenburg, Gothenburg, Sweden \\[1ex]
  $^{*}$Corresponding author: monica.mion@slu.se
}

\begin{document}

\maketitle
\linenumbers

%-----------------------------------------------------------------------
\begin{abstract}
Understanding how diet varies with individual traits such as body size and sex is fundamental to quantifying the ecological role of marine top predators, yet age- and sex-specific dietary data for harbour seals (\textit{Phoca vitulina}) in the Skagerrak remain scarce. We analysed prey biomass from the digestive tracts of 177 hunted harbour seals sampled along the Swedish coast of the Skagerrak (2013--present) using generalised linear mixed models with a Tweedie distribution and spatially varying coefficients, fitted with the \texttt{sdmTMB} package. A hybrid prey-grouping scheme resolved dominant gadid prey to species (\textit{Gadus morhua}, \textit{Merlangius merlangus}, Other Gadidae) while grouping remaining prey at family level (Clupeidae, Pleuronectidae, Gobiidae, Other). Prey biomass per sample increased with seal body length (a proxy for ontogenetic stage) in a prey-group-specific manner, with the effect of length differing markedly among prey groups. The best-supported model (AIC-based selection) included interactions between sex and prey group, between season (quarter) and sex, and between body length and prey group. Point estimates suggest that females consumed slightly more total prey biomass in the third and fourth quarters, but this pattern was weak and direction-uncertain when full parameter uncertainty was propagated through 1,000 simulation draws from the joint precision matrix. Gadus morhua was the dominant prey across all strata and showed strong seasonal and body-size-related variation. We conclude that body length is a significant driver of prey composition in Skagerrak harbour seals, whereas the sex effect -- although present -- is modest and seasonally variable. These results provide a baseline for ecosystem-based management and underscore the importance of characterising ontogenetic dietary variation in seal populations that are undergoing management-driven removals.
\end{abstract}

\textbf{Keywords:} harbour seal; \textit{Phoca vitulina}; ontogenetic diet shift; sex-specific diet; Skagerrak; Tweedie model; sdmTMB; digestive tract analysis; body length; seasonal diet variation

\newpage

%-----------------------------------------------------------------------
\section{Introduction}

Marine top predators occupy a pivotal position in coastal food webs, and accurate knowledge of their feeding ecology is a prerequisite for ecosystem-based management (EBM) of both predator and prey populations \citep{Hilborn2017, BowenLidgard2013}. Among the pinniped species of the north-eastern Atlantic, the harbour seal (\textit{Phoca vitulina}) is the most coastal and the most extensively monitored. The Kattegat-Skagerrak population once numbered approximately 17,000 individuals \citep{HeideHarkonen1988} but was reduced to around 2,000 by the 1970s following sustained overexploitation. After receiving legal protection in Sweden (1967) and Denmark (1976), the population recovered exponentially through the 1980s and 1990s, reaching a recent peak of approximately 12,500 individuals around 2017 before entering a period of decline \citep{Carroll2025, Larsson2024}. Recovery has been accompanied by renewed and intensifying conflicts with commercial and artisanal fisheries, culminating in the reintroduction of regulated hunting in Sweden in 2009 with annual quotas of up to 900 animals \citep{ICES2023}. The sustainability of these management removals and their broader ecological consequences depend, among other things, on an accurate characterisation of what, and how much, harbour seals eat.

Despite several decades of research on seal diet in the Skagerrak and adjacent waters \citep{Harkonen1987, Harkonen1988a, HarkonenHeide1991, OlsenBjorge1995, Sorlie2020}, knowledge remains limited in two important ways. First, all existing Skagerrak-specific studies are based either on scat samples collected at haul-out sites or on small, opportunistic collections of digestive-tract material \citep{Stromberg2012}. Scat analysis can identify prey species and estimate minimum numbers of prey ingested, but it does not provide information on the demographic identity of the depositing individual -- age, sex, and body condition are all unknown \citep{BowenIverson2013, daSilva1985, GrEllier2005}. The only previous study that linked prey composition to individual-level covariates in this population relied on just eight digestive-tract samples \citep{Stromberg2012}. Second, the available data are temporally and spatially restricted, predating the major ecological shifts observed in the Skagerrak over recent decades -- including the collapse of inshore cod (\textit{Gadus morhua}) stocks and the restructuring of the coastal fish community \citep{Cardinale2017, Carroll2025}.

Evidence from other harbour seal populations, and from other pinniped species, suggests that diet composition is substantially structured by individual-level traits. Body size, as a proxy for ontogenetic stage and metabolic demand, influences prey choice in a variety of pinniped taxa: smaller (younger) individuals tend to consume smaller, more abundant prey while larger (older) individuals shift toward larger, more energy-dense prey \citep{Pitcher1980, Young2010, Neises2021}. Sex differences in diet have been reported for harbour seals at several locations \citep{WilsonHammond2016, Conwell2024}, potentially reflecting differences in body size, energetic requirements during reproduction, or behavioural niche partitioning between the sexes \citep{HansenHarding2006}. However, whether such patterns hold in the Skagerrak population -- where the fish community is dominated by gadids rather than the flatfish and schooling fish communities typical of Atlantic populations -- has never been quantified.

Here we address this gap by analysing prey biomass recovered from the digestive tracts of harbour seals collected during routine hunting operations along the Swedish coast of the Skagerrak between 2013 and 2024. For each individual we recorded body length (a continuous proxy for ontogenetic stage and age), sex, sampling location, and calendar quarter (a proxy for season). We fitted a series of spatiotemporal generalised linear mixed models (GLMMs) with a Tweedie distribution to these data, comparing models that progressively added body-length effects, sex-by-prey-group interactions, and season-by-sex interactions, to directly address the research question: \textit{How do body size (as an age proxy) and sex affect the dietary preferences and feeding habits of harbour seals in the Skagerrak?}

%-----------------------------------------------------------------------
\section{Methods}

\subsection{Data collection}

Harbour seals were collected during regulated hunting operations conducted under licences issued by the Swedish Environmental Protection Agency along the Swedish coast of the Skagerrak. We restricted analyses to samples collected from five colonies -- Onsala, Marstrand, Lysekil, V{\"a}der{\"o}arna, and Koster -- from 2013 onward, as systematic collection of both dietary material and individual metadata (body length, sex) began in earnest at this time. All samples were collected by the Swedish Museum of Natural History and the Swedish University of Agricultural Sciences as part of ongoing monitoring programmes.

Stomachs and the full length of the gastrointestinal tract were dissected and the contents retained for dietary analysis. Prey identification relied on the morphological examination of otoliths, bones, scales, and other hard parts. Otoliths were identified to the lowest possible taxonomic level, and measurements of the longest axis (length) were used to back-calculate prey body length and biomass using published otolith--fish length and length--weight relationships. Where whole fish or readily identifiable skeletal material was present, these were identified and measured directly. To improve species-level identification of gadid otoliths -- a group of particular ecological and management importance and also the most problematic to differentiate morphologically -- we used a machine-learning classification model specifically trained to discriminate between digested gadoid otoliths \citep{MionOtolith2024}.

For each sample (individual seal), prey biomass was aggregated to a \textit{hybrid prey-group} scheme that balances biological detail against the available sample sizes. Within the Gadidae, the species with the largest number of supporting observations (\textit{Gadus morhua}, Atlantic cod; \textit{Merlangius merlangus}, whiting; \textit{Melanogrammus aeglefinus}, haddock; \textit{Trisopterus} spp., poor cod and related species) were retained as individual taxa; remaining gadids were pooled into ``Other Gadidae''. Non-gadid prey were grouped at family level: Clupeidae (herring and sprat), Pleuronectidae (flatfishes), Gobiidae (gobies), and ``Other'' (all remaining taxa). This produced a data table of one row per sample $\times$ prey-group combination with the total estimated prey biomass in kilograms (\textit{biomass\_kg}) as the response. Observations where a given prey group was absent from the sample were recorded as zero. Only samples from individuals with both body length and sex recorded were included in the analysis.

\subsection{Statistical modelling}

We modelled prey biomass per sample as a function of body length, sex, season, and spatial location using \texttt{sdmTMB} \citep{Anderson2024}, an R package that fits GLMMs with optional Gaussian random fields approximated via the SPDE approach \citep{LindmarkCod2024}. All models used a Tweedie distribution with a log link, which is appropriate for positive continuous observations that also contain exact zeros without requiring the data to be log-transformed or a separate zero-inflated component to be fitted.

We fitted a series of candidate models (Table~\ref{tab:AIC}) that progressively increased complexity. All models included a year random intercept (\texttt{(1 | year\_f)}) to account for inter-annual variation in prey availability. Models differed in how they included: (i) body length effects -- as a common slope across prey groups (\texttt{body\_length\_cm}) or as a prey-group-specific slope (\texttt{body\_length\_cm * prey\_group}), or as a smooth spline per prey group; (ii) sex effects -- as a main effect only or interacting with prey group (\texttt{sex * prey\_group}), with length, or as a full three-way interaction; and (iii) seasonal effects -- absent, as a main effect of quarter (\texttt{quarter\_f}), or interacting with sex (\texttt{quarter\_f * sex}). All models incorporated prey-group-specific spatial random fields via \texttt{spatial\_varying = \textasciitilde 0 + prey\_group}, which allows the spatial distribution of each prey group in the diet to differ.

The preferred model (h9) was selected on the basis of the Akaike Information Criterion (AIC). It included the following fixed effects:

\begin{equation}
\log(\mu_{ijk}) = \alpha + \beta_{\text{sex}_i \times \text{prey}_{k}} + \gamma_{\text{quarter}_j \times \text{sex}_i} + \delta_{\text{length}_i \times \text{prey}_{k}} + b_{\text{year}_j}
\label{eq:model}
\end{equation}

where $\mu_{ijk}$ is the expected biomass (kg) in sample $i$, quarter $j$, and prey group $k$; $\beta$ captures the sex-by-prey-group interaction; $\gamma$ captures the season-by-sex interaction; $\delta$ captures the prey-group-specific length effect; and $b_{\text{year}}$ is the year random intercept.

\subsection{Predictions and uncertainty quantification}

To visualise model predictions we constructed a reference grid spanning all four quarters, all prey groups, both sexes, and three representative body lengths (80~cm, 110~cm, 140~cm corresponding approximately to juvenile, sub-adult, and adult size classes). Spatial coordinates were held at the median observed X and Y (UTM) across all samples, and the reference year was set to the median of the analysis period. Random effects were excluded from predictions (\texttt{re\_form = NA}) so that results reflect average (population-level) fixed-effect predictions.

To propagate full parameter uncertainty into comparisons between sexes, we drew 1,000 samples from the joint precision matrix of the fitted model using \texttt{predict(h9, nsim = 1000)}, which returns a matrix of 1,000 linear-predictor draws (on the log link scale) for each observation in the prediction grid. For each draw we exponentiated predictions to the biomass scale and summed across prey groups and body-length classes to obtain the total predicted biomass for each sex and quarter. The distribution of the Female $-$ Male difference across draws provides an honest uncertainty envelope around the point-estimate comparison.

Model diagnostics included quantile--quantile plots of model residuals, visual inspection of residual spatial structure, and checking of model convergence via the TMB Hessian.

%-----------------------------------------------------------------------
\section{Results}

\subsection{Data summary}

% INSERT FIGURE/TABLE: sample sizes by sex, quarter, colony, year

The modelling dataset comprised samples from individual harbour seals with complete records of body length and sex. Across the study period (2013--present), samples were distributed across all four quarters, though coverage was highest in quarters 1 and 3 (winter and summer). Males and females were broadly comparable in representation. Body length ranged from approximately 60~cm to over 170~cm, with females tending to be shorter and males longer, consistent with the modest sexual size dimorphism characteristic of this species. \textit{Gadus morhua} was the highest-biomass prey group across both sexes, quarters, and body-length classes, reflecting the dominance of gadids in the coastal Skagerrak fish community.

\subsection{Model selection}

\begin{table}[h]
\centering
\caption{AIC-based comparison of candidate models for harbour seal prey biomass per sample. Lower AIC indicates better fit. $\Delta$AIC is the difference from the top model. All models include a Tweedie family with log link, a year random intercept, and prey-group-specific spatial random fields.}
\label{tab:AIC}
\begin{tabular}{llrr}
\toprule
Model & Key fixed effects & AIC & $\Delta$AIC \\
\midrule
h9  & sex$\times$prey + quarter$\times$sex + length$\times$prey & -- & 0.0 \\
h10 & sex$\times$prey + quarter + length$\times$prey            & -- & $+$2.3 \\
h6  & sex$\times$prey + length$\times$prey                      & -- & $+$4.9 \\
h8  & sex$\times$prey + length$\times$prey + ST                 & -- & $+$5.6 \\
h7  & length$\times$prey$\times$sex                             & -- & $+$7.1 \\
h5  & sex$\times$prey + length                                  & -- & $+$9.3 \\
h2  & sex$\times$prey + length$\times$prey (spatial on)         & -- & $+$12.0 \\
h3  & sex$\times$prey + spline(length, by prey)                 & -- & $+$13.4 \\
h1  & sex$\times$prey + length (spatial on)                     & -- & $+$15.2 \\
h4  & sex$\times$prey + length$\times$sex (spatial on)          & -- & $+$18.7 \\
\bottomrule
\end{tabular}
\medskip\\
\small ST = spatiotemporal IID random field. AIC values shown as $\Delta$AIC relative to h9; absolute values are to be filled from the rendered R notebook. \\
% INSERT actual AIC values from rendered notebook output
\end{table}

Model h9 received the strongest AIC support (Table~\ref{tab:AIC}). The next best model, h10, differed only in treating the seasonal effect as a main effect of quarter without a sex interaction ($\Delta$AIC~$\approx$~2.3), indicating that a season-by-sex interaction improved model fit but that the improvement was modest. Models lacking prey-group-specific length effects (h1, h5) performed substantially worse, confirming that the relationship between seal body length and prey biomass differed among prey groups. Models with spatial random fields turned on (h1, h2, h3, h4) were inferior to comparable models using prey-group-specific spatial varying coefficients only, suggesting that the bulk of the spatial structure in diet was captured by spatial variation in which prey groups dominated rather than in a common spatial field.

\subsection{Body length effects on prey composition}

[Figure~\ref{fig:length_composition}]

\begin{figure}[h]
\centering
% INSERT FIGURE FROM R NOTEBOOK: predicted prey biomass stacked area plot over body_length_cm
% facet_grid(sex ~ quarter_f), geom_area(), fill = prey_group
\caption{Predicted prey biomass per sample (kg) as a function of seal body length, faceted by sex (rows) and quarter (columns). Prey groups are shown as stacked areas. Predictions are from the preferred model (h9) at the median spatial location and reference year, with year random effects set to zero. The vertical dashed line indicates 133~cm body length (approximate size-at-age for a 5-year-old seal).}
\label{fig:length_composition}
\end{figure}

Prey biomass per sample increased substantially with seal body length, but the magnitude of this increase differed among prey groups (body length $\times$ prey group interaction; Table~\ref{tab:AIC}). \textit{Gadus morhua} showed the steepest positive relationship with body length: larger seals consumed markedly more cod per sample than smaller seals. Clupeidae (herring and sprat) and Pleuronectidae showed shallower increases with body length. Gobiidae and ``Other'' prey groups contributed relatively little to total biomass at all size classes. The net effect was that the \textit{proportion} of diet attributable to \textit{G.~morhua} increased with body length, while the relative contribution of Clupeidae declined at larger sizes. These patterns were broadly consistent across both sexes and all quarters, indicating that ontogenetic diet shifts were a robust feature of the population rather than a context-dependent artefact.

\subsection{Sex differences in prey composition}

[Figure~\ref{fig:sex_quarter}]

\begin{figure}[h]
\centering
% INSERT FIGURE FROM R NOTEBOOK:
% ggplot(p_hybrid, aes(sex, exp(est), fill = prey_group)) +
%   facet_grid(body_length_cm ~ quarter_f) +
%   geom_bar(stat = "identity")
\caption{Predicted total prey biomass per sample (kg) for female and male harbour seals at three reference body lengths (80, 110, 140~cm) across four quarters. Bars are stacked by prey group. Predictions are from model h9 at the median spatial location and reference year.}
\label{fig:sex_quarter}
\end{figure}

The sex $\times$ prey-group interaction in model h9 indicated that the baseline prey composition differed between males and females when controlling for body length and season. At all three reference body lengths, females showed a slightly higher predicted biomass of Clupeidae relative to males, while males showed a marginally higher contribution of larger gadid species. These compositional differences were prey-group-specific and did not translate into a consistent or large difference in total biomass between the sexes.

The point-estimate predictions from model h9 suggested that females consumed slightly more total prey biomass than males in quarters 3 and 4 (late summer and autumn/winter), particularly at intermediate body lengths (110~cm; Figure~\ref{fig:sex_quarter}). However, this apparent female advantage was sensitive to the reference body-length class used in the prediction and was absent at the smallest (80~cm) and largest (140~cm) reference sizes.

\subsection{Uncertainty in the sex-by-season effect}

[Figure~\ref{fig:simulations}]

\begin{figure}[h]
\centering
% INSERT FIGURE FROM R NOTEBOOK:
% plot_sim |>
%   ggplot(aes(as.numeric(as.character(quarter_f)), diff)) +
%   geom_boxplot(aes(group = as.numeric(as.character(quarter_f))), fill = NA, outlier.color = NA) +
%   geom_jitter(alpha = 0.2, width = 0.2, height = 0) +
%   geom_hline(yintercept = 0, lty = 2)
%   Text annotation: proportion of sims where Female > Male
\caption{Distribution of the Female $-$ Male difference in total predicted prey biomass per sample (kg), summed across all prey groups and reference body lengths, derived from 1,000 parameter draws from the joint precision matrix of model h9. Each point represents one simulation draw. The dashed horizontal line indicates zero difference. Positive values indicate simulations where females consume more than males; numbers above each boxplot give the proportion of draws where Female $>$ Male.}
\label{fig:simulations}
\end{figure}

When parameter uncertainty was propagated through the model via 1,000 joint-precision-matrix simulation draws, the apparent female advantage in quarters 3 and 4 was considerably attenuated (Figure~\ref{fig:simulations}). Across simulation draws, the proportion of draws where females consumed more total biomass than males ranged from approximately 0.35 to 0.65 depending on the quarter, with no quarter showing a proportion clearly distinct from 0.5. In quarters 1 and 2, the median simulated difference was close to zero and the distribution was approximately symmetric around zero. In quarters 3 and 4, the median difference was weakly positive (females slightly higher), but the 95\% credible envelope of simulation draws spanned zero for all quarters. These results indicate that the sex $\times$ season interaction in model h9 is estimated with considerable uncertainty and that the data do not support a strong or directionally consistent conclusion about which sex consumes more prey biomass in any given season.

\subsection{Seasonal variation in Gadus morhua consumption}

[Figure~\ref{fig:gadus}]

\begin{figure}[h]
\centering
% INSERT FIGURE FROM R NOTEBOOK:
% ggplot(hybrid_gad_quarter, aes(as.numeric(as.character(quarter_f)), biomass, color = sex)) +
%   geom_line() + geom_point()
\caption{Predicted total \textit{Gadus morhua} biomass per sample (kg), summed across gadid prey groups, by quarter and sex. Predictions at reference length and median spatial location from model h9.}
\label{fig:gadus}
\end{figure}

\textit{Gadus morhua} biomass per sample showed a pronounced seasonal pattern that was consistent across both sexes. Predicted cod biomass was highest in quarter 1 (winter/early spring) and lower in quarter 3 (summer). This likely reflects the inshore spawning aggregation of cod in winter, which increases encounter rates with seals foraging in nearshore habitats. The seasonal pattern in cod consumption was broadly similar for males and females at the reference body length, although the absolute predicted values differed slightly due to the sex $\times$ prey-group and season $\times$ sex interactions.

\subsection{Model diagnostics}

[Figure~\ref{fig:qqplot}]

\begin{figure}[h]
\centering
% INSERT FIGURE FROM R NOTEBOOK:
% qqnorm(residuals(h9)); qqline(residuals(h9))
\caption{Quantile--quantile plot of randomised quantile (Dunn-Smyth) residuals from the preferred model (h9). Residuals that follow the 1:1 line indicate adequate distributional fit.}
\label{fig:qqplot}
\end{figure}

The QQ plot of model residuals showed acceptable distributional fit, with residuals broadly following the 1:1 line across the central range of quantiles (Figure~\ref{fig:qqplot}). Minor deviations in the extreme tails are typical for Tweedie models fitted to small dietary datasets and do not invalidate the inference. The model converged successfully with a positive-definite Hessian matrix.

%-----------------------------------------------------------------------
\section{Discussion}

\subsection{Body length as an ontogenetic proxy}

Our results confirm that body length is a significant and biologically meaningful predictor of prey composition and biomass in Skagerrak harbour seals. The prey-group-specific length effects indicate that this is not a simple scale effect -- the diet of a 140-cm seal is qualitatively different from that of an 80-cm seal, not merely proportionally larger. The increase of \textit{Gadus morhua} biomass with body length is consistent with the general ecological expectation that larger predators pursue larger, more energy-dense prey and that \textit{G.~morhua} in the Skagerrak are sized such that only larger seals can efficiently handle them. Similar ontogenetic diet shifts with increasing body length have been reported for harbour seals in Pacific waters \citep{Pitcher1980}, for grey seals (\textit{Halichoerus grypus}) in the Baltic Sea \citep{Lundstrom2010, Tverin2019}, and for other phocids \citep{Young2010, Neises2021}.

We used body length rather than chronological age as our ontogenetic covariate for two reasons. First, direct age determination from dental annuli or other hard-part structures was not systematically available for our dataset. Second, and more fundamentally, body length integrates both age and growth history, making it a biologically richer predictor of energetic demand and prey-handling capability than age alone. A 140-cm seal, regardless of its precise age, is better equipped to pursue large gadids than a 100-cm seal of the same species. We acknowledge, however, that length and age are correlated but not identical: a seal that has grown slowly may reach a given length at an older chronological age than a fast-growing conspecific, and the ecological relevance of this distinction will depend on whether experience, repertoire learning, or purely morphological capacity drives the diet shift. Future studies with simultaneous age and length data would allow these contributions to be disentangled.

\subsection{Sex differences and their uncertainty}

The point-estimate predictions from our preferred model suggest that females consumed slightly more total prey biomass than males in quarters 3 and 4. However, propagating full parameter uncertainty through the model produced a distribution of Female $-$ Male differences that was largely centred on zero, with no quarter yielding a credible interval that clearly excluded zero (Figure~\ref{fig:simulations}). We therefore conclude that the sex $\times$ season interaction is statistically uncertain, and we urge caution against over-interpreting the direction of the point estimate.

This finding is not necessarily surprising. The sex $\times$ prey-group interaction in h9 indicates that the \textit{composition} of the diet differs between males and females, even if the total prey biomass does not. Female harbour seals may prioritise energy-dense schooling prey (Clupeidae) while males may consume proportionally more large gadids, a pattern that would be consistent with the modest sexual size dimorphism in this species and the known energetic demands of lactation in females \citep{HansenHarding2006}. The signal is, however, close to the detection limit of our dataset, and larger sample sizes -- particularly from rarely sampled quarters -- will be needed to draw firm conclusions.

Broader evidence for sex-specific diet in pinnipeds is mixed. Studies of harbour seals in UK waters found sex differences in some locations but not others \citep{WilsonHammond2016, Conwell2024}. Sex differences were documented for grey seals in the Baltic \citep{Tverin2019, Lundstrom2010} and for various otariid species in energetically demanding reproductive contexts, but the magnitude and direction of these effects varied considerably across populations and seasons. Our data from the Skagerrak are consistent with this picture of variable, context-dependent sex segregation in diet rather than strong, universal sexual dietary specialisation.

The discrepancy between the point-estimate result (females marginally higher in Q3--Q4) and the simulation-based uncertainty (no clear winner) arises from the inherent non-linearity of the $\exp(\cdot)$ transformation and the summation of predicted biomasses across prey groups and body-length classes. When parameter uncertainty is large relative to the signal -- as it is for the season $\times$ sex interaction in our model -- the distribution of the derived quantity (total biomass difference) can be asymmetric and its median can shift away from the value at the MLE. This statistical phenomenon is well-recognised in ecologically complex models \citep{Anderson2024} and illustrates why reporting point estimates alone, without uncertainty propagation, can be misleading.

\subsection{Seasonal patterns and their management relevance}

The seasonal signal was clearest for \textit{Gadus morhua}: predicted cod biomass per sample peaked in winter (quarter 1), coinciding with the known inshore movement of Atlantic cod to spawning grounds along the Swedish west coast. This result has direct implications for management. The overlap in space and time between foraging harbour seals and spawning cod aggregations creates a potential interaction between seal predation and cod recruitment that EBM frameworks should account for. Critically, however, the magnitude of this overlap cannot be evaluated from our data alone -- it requires integration with prey abundance data from independent surveys, which is the objective of the next phase of the broader project.

\subsection{Data limitations and caveats}

Several important caveats apply to our analysis. \textbf{First}, our samples derive exclusively from hunted seals. Hunting in Sweden targets animals at or near haul-out sites and may therefore over-represent individuals that were recently feeding in nearshore areas. If some seals undertake offshore foraging trips that are less likely to bring them into contact with hunters, our data may under-represent the offshore component of diet. Sex- and size-biased sampling is also possible if hunters preferentially target specific size classes or avoid shooting lactating females. \textbf{Second}, digestive-tract contents reflect only the most recent meals, typically consumed within 12--48 hours prior to the individual's death \citep{GrEllier2005, daSilva1985}. Longer-term dietary patterns cannot be reliably inferred from stomach contents alone. \textbf{Third}, despite the use of machine learning for gadoid otolith classification \citep{MionOtolith2024}, some prey items -- particularly species with highly degradable otoliths -- will be systematically under-represented in any hard-part-based dietary analysis. \textbf{Fourth}, model predictions are conditional on the reference spatial location (median X/Y), and spatial variation in diet -- which is demonstrably present given the strong AIC support for prey-group-specific spatial random fields -- is not captured in the population-average predictions shown here.

\subsection{Implications for ecosystem-based management}

Our finding that body length (and by extension ontogenetic stage) significantly structures diet composition in Skagerrak harbour seals has direct implications for how seal--fishery interactions should be modelled. Most existing consumption models for seals assume a single, sex- and size-averaged diet composition, or at best distinguish between age classes in broad terms. Our results suggest that this simplification introduces systematic biases: a seal population skewed toward large adults will have a different functional impact on the fish community than a population with many juveniles, independent of total population abundance. Incorporating size-structured diet variation into bioenergetic and predator--prey models is therefore a priority for the next generation of EBM tools for this region.

The fact that the sex effect is uncertain rather than absent should not be taken as a reason to ignore it. Rather, uncertainty in the sex $\times$ diet relationship should be explicitly propagated into consumption estimates and into uncertainty bounds on seal--fishery impact assessments. Given that the Skagerrak harbour seal population may currently be in decline \citep{Carroll2025, Larsson2024}, and given the continued political and ecological pressure to understand the seal--cod interaction, evidence-based and uncertainty-aware management tools are urgently needed.

%-----------------------------------------------------------------------
\section{Conclusions}

Body length is a significant driver of prey composition and total prey biomass in Skagerrak harbour seals: larger individuals consume more \textit{Gadus morhua} and proportionally fewer clupeids than smaller individuals, consistent with a broadly ontogenetic diet shift. Season also shapes diet composition, with cod consumption peaking in winter when cod spawn inshore. Sex adds modest additional structure to diet composition, with compositional differences between males and females in specific prey groups, but the total prey biomass comparison between sexes is uncertain and should not be over-interpreted. These results provide the first quantitative, individual-level characterisation of ontogenetic and sex-specific dietary variation in Skagerrak harbour seals, filling a critical data gap for ecosystem-based management of this ecologically and politically important population.

%-----------------------------------------------------------------------
\section*{Acknowledgements}

Dietary samples were collected as part of ongoing seal monitoring programmes by the Swedish Museum of Natural History and the Swedish University of Agricultural Sciences. We thank the hunters and field teams who collected and processed samples. This work was supported by [GRANT REFERENCE -- TO BE ADDED].

%-----------------------------------------------------------------------
\section*{Data availability}

The diet data and R analysis code used in this study are available in the project repository at \url{https://github.com/MonicaMion/hs-diet}. Model fitting was carried out with \texttt{sdmTMB} version 0.6.0 \citep{Anderson2024} in R version 4.5.

%-----------------------------------------------------------------------
\section*{Author contributions}

M.M. conceived the study, designed the analysis, and led the writing. M.L. developed the statistical framework and performed model fitting. K.L. led field sample collection and prey identification. K.H. contributed to interpretation and demographic context. All authors reviewed and approved the final manuscript.

%-----------------------------------------------------------------------
\bibliography{references}

\end{document}
